\documentclass[11pt]{article}
\usepackage[margin=24mm]{geometry}
\usepackage{amsmath,amssymb,amsthm,amscd}
\usepackage[T1]{fontenc}
\usepackage{lmodern}
\usepackage[colorlinks=true,urlcolor=blue,linkcolor=blue,hypertexnames=false]{hyperref}
\newcommand{\un}{\mathrm{un}}
\newcommand{\al}{\mathrm{al}}
\newcommand{\Gal}{\operatorname{Gal}}
\newcommand{\End}{\operatorname{End}}
\newcommand{\Frob}{\operatorname{Frob}}
\newcommand{\m}{\mathfrak m}
\newcommand{\n}{\mathfrak n}
\newcommand{\terms}[1]{\text{terms of degree }\geq #1}
\newcommand{\Khat}{\widehat K^{\un}}
\newtheorem*{target}{Theorem I.3.9}
\newtheorem*{intertwiner}{Lemma I.2.11 (support)}
\newtheorem*{comparison}{Proposition I.3.10 (support)}
\newtheorem*{difference}{Lemma I.3.11 (support)}
\newtheorem*{closedfield}{Lemma I.3.12 (support)}
\title{Independence of the Lubin--Tate parameter\\and the local Artin action}
\author{J. S. Milne\\\small Source: \emph{Class Field Theory}, v4.03, 6 August 2020}
\date{}
\begin{document}
\maketitle
\noindent\small
Private reading transcription of the author's Theorem I.3.9 and its supporting argument.
The original PDF is retained at
\url{https://www.jmilne.org/math/CourseNotes/CFT.pdf}.
The main argument is on physical PDF pages 49--52 (printed pages 40--43).
Context below comes from physical pages 41--42, 44--45 and 47.
Original numbering is prefixed by the chapter to disambiguate it.
Mathematical expressions and the author's abbreviated steps are retained;
connective prose and layout are normalized. This is published-material transcription,
not author-native TeX or a new proof. Author errata are identified separately below.
\normalsize

\section*{Author context and notation}
Let $K$ be a nonarchimedean local field, $A=\mathcal O_K$, with maximal ideal
$\m$, residue field $k=A/\m$ of order $q$, and prime element $\pi$.
Write $K^{\al}$ for the separable algebraic closure used in the source and
$K^{\un}$ for its maximal unramified subextension.
The source's class $\mathcal F_\pi$ consists of series $f\in A[[T]]$ satisfying
\[
 f(T)=\pi T+\terms{2},\qquad f(T)\equiv T^q\pmod\pi.
\]
Summary I.2.20 provides the unique formal group law $F_f$ admitting $f$ as an
endomorphism and its unique $A$-module structure $a\mapsto[a]_f$, with
\[
 [a]_f=aT+\terms{2},\qquad [a]_f\circ f=f\circ[a]_f,\qquad [\pi]_f=f.
\]
For $g\in\mathcal F_\pi$, $F_f$ and $F_g$ have the canonical $A$-isomorphism
denoted $[1]_{g,f}:F_f\longrightarrow F_g$ in the source.

The setup in Section I.3 is
\[
 \Lambda_f=\{\alpha\in K^{\al}:|\alpha|<1\},\quad
 \alpha+_{\Lambda_f}\beta=F_f(\alpha,\beta),\quad
 a*\alpha=[a]_f(\alpha).
\]
Let $\Lambda_n$ be the submodule annihilated by $[\pi]_f^n$; equivalently,
its elements are the roots of $f^{(n)}$ of absolute value less than one.
Set $K_{\pi,n}=K[\Lambda_n]$ and $K_\pi=\bigcup_nK_{\pi,n}$.
The relevant preceding conclusions of I.3.4 and I.3.6 are
\[
 \Lambda_n\simeq A/(\pi^n),\quad
 [K_{\pi,n}:K]=(q-1)q^{n-1},\quad
 (A/\m^n)^\times\xrightarrow{\sim}\Gal(K_{\pi,n}/K),
\]
where the last map is induced by the $A$-action, and $K_{\pi,n}/K$ is totally ramified.
These preceding Lubin--Tate structure results are context, not additional counted targets.
The author's separate errata include the convergence convention for $\Lambda_f$;
the original setup is preserved here and that note remains attached to the source package.

% BEGIN SUPPORT LEMMA I.2.11
\section*{The author's intertwining lemma}
\begin{intertwiner}
Let $f,g\in\mathcal F_\pi$, and let $\phi_1(X_1,\ldots,X_n)$ be a linear form
with coefficients in $A$. There is a unique $\phi\in A[[X_1,\ldots,X_n]]$ such that
\[
 \begin{cases}
 \phi(X_1,\ldots,X_n)=\phi_1+\terms{2},\\
 f(\phi(X_1,\ldots,X_n))=\phi(g(X_1),\ldots,g(X_n)).
 \end{cases}
\]
\end{intertwiner}
\begin{proof}[Author proof]
Inductively construct the unique polynomial $\phi_r$ of degree $r$ with
\[
 \begin{cases}
 \phi_r(X_1,\ldots,X_n)=\phi_1+\terms{2},\\
 f(\phi_r(X_1,\ldots,X_n))
 =\phi_r(g(X_1),\ldots,g(X_n))+\terms{r+1}.
 \end{cases}
\]
For $r=1$ the candidate is $\phi_1$. If $\phi_1=\sum a_iX_i$, its second
condition is $\pi(\sum a_iX_i)=\sum a_i(\pi X_i)+\terms{2}$.
After $\phi_r$ has been constructed, write $\phi_{r+1}=\phi_r+Q$, with $Q$
homogeneous of degree $r+1$. The two sides of the required relation are
\begin{align*}
 f(\phi_{r+1}(X_1,\ldots,X_n))
 &=f(\phi_r(X_1,\ldots,X_n))+\pi Q(X_1,\ldots,X_n)+\terms{r+2},\\
 \phi_{r+1}(g(X_1),\ldots,g(X_n))
 &=\phi_r(g(X_1),\ldots,g(X_n))+Q(\pi X_1,\ldots,\pi X_n)+\terms{r+2}.
\end{align*}
The source's homogeneity line prints
$Q(\pi X_1,\ldots)=\pi^{r+1}Q(\pi X_1,\ldots)$;
the author's erratum for printed p.32 deletes the third $\pi$, giving
$\pi^{r+1}Q(X_1,\ldots)$.
The following displayed requirement in the source is
\[
 (\pi^{r+1}-\pi)Q(X_1,\ldots,X_n)
 =f(\phi_r(X_1,\ldots,X_n))-\phi_r(g(X_1),\ldots,g(X_n))
 +\terms{r+2}.
\]
Thus $Q$ is determined by
\[
 \frac{f(\phi_r(X_1,\ldots,X_n))-\phi_r(g(X_1),\ldots,g(X_n))}
 {(\pi^r-1)\pi}
 =Q+\terms{r+2}.
\]
The characteristic-$p$ binomial calculation gives
\[
 f\circ\phi_r-\phi_r\circ g
 \equiv\phi_r(X_1,\ldots,X_n)^q-\phi_r(X_1^q,\ldots,X_n^q)
 \equiv0\pmod\pi.
\]
The numerator is divisible by $\pi$, and $\pi^r-1$ is a unit of $A$.
Consequently $Q$ has coefficients in $A$ and, by the induction hypothesis,
has degree $r+1$. The compatible polynomials determine a unique series with
\[
 \phi_{r+1}=\phi_r+\terms{r+1},\qquad
 \phi=\phi_r+\terms{r+1}\quad\text{for every }r.
\]
The source ends the verification with
\begin{align*}
 f(\phi(X_1,\ldots,X_n))
 &=f(\phi_r(X_1,\ldots,X_n))+\terms{r+1}\\
 &=\phi_r(g(X_1),\ldots,g(X_n))+\terms{r+1}\\
 &=\phi(f(X_1),\ldots,f(X_n))+\terms{r+1}.
\end{align*}
The last line is transcribed with $f$ as printed in the original PDF.
The author concludes that, as this holds for all $r$, the second required property holds.
\end{proof}
% END SUPPORT LEMMA I.2.11

% BEGIN PRIMARY STATEMENT I.3.9
\section*{The local Artin map and the target statement}
Define
\[
 \phi_\pi:K^\times\longrightarrow\Gal((K_\pi\cdot K^{\un})/K).
\]
Because $K_\pi\cap K^{\un}=K$, the two restrictions specify its action.
For $a=u\pi^m$, $u\in A^\times$, its restriction to $K^{\un}$ is $\Frob_K^m$,
and its action on the torsion generators of $K_\pi$ is
\[
 \phi_\pi(a)(\lambda)=[u^{-1}]_f(\lambda),\qquad
 \lambda\in\bigcup_n\Lambda_n.
\]
\begin{target}
Both $K_\pi\cdot K^{\un}$ and $\phi_\pi$ are independent of the choice of $\pi$.
\end{target}
% END PRIMARY STATEMENT I.3.9

Recall that $K^{\un}$ is not complete. Let $\Khat$ be its completion, and let
$B$ be its valuation ring. Write $\sigma$ for $\Frob_K$ and its extension to
$\Khat$. For $\theta(T)=\sum b_iT^i\in B[[T]]$, set
$(\sigma\theta)(T)=\sum(\sigma b_i)T^i$.

% BEGIN SUPPORT PROPOSITION I.3.10 STATEMENT
\begin{comparison}
Let $f\in\mathcal F_\pi$ and $g\in\mathcal F_\varpi$, where $\varpi=u\pi$.
The formal group laws $F_f,F_g$ become $A$-isomorphic over $B$.
More precisely, there are $\varepsilon\in B^\times$ and $\theta(T)\in B[[T]]$
with $\sigma\varepsilon=\varepsilon u$ such that
\begin{align*}
 \text{(a)}\quad&\theta(T)=\varepsilon T+\terms{2},\\
 \text{(b)}\quad&\sigma\theta=\theta\circ[u]_f,\\
 \text{(c)}\quad&\theta(F_f(X,Y))=F_g(\theta(X),\theta(Y)),\\
 \text{(d)}\quad&\theta\circ[a]_f=[a]_g\circ\theta.
\end{align*}
\end{comparison}
Conditions (c),(d) express compatibility with the formal groups and $A$-actions;
(a) makes the homomorphism an isomorphism, since $\varepsilon$ is a unit.
% END SUPPORT PROPOSITION I.3.10 STATEMENT

% BEGIN SUPPORT LEMMA I.3.11
\begin{difference}
The maps
\[
 b\longmapsto\sigma b-b:B\longrightarrow B,
 \qquad b\longmapsto\frac{\sigma b}{b}:B^\times\longrightarrow B^\times
\]
are surjective, with kernels $A$ and $A^\times$, respectively.
\end{difference}
\begin{proof}[Author proof]
Let $R$ be the valuation ring of $K^{\un}$ and $\n$ its maximal ideal.
Then $R$ is a discrete valuation ring and $\varprojlim R/\n^n=B$ (I.A.7).
The author proves inductively that
\begin{equation}
 0\longrightarrow A/\m_K^n\longrightarrow R/\n^n
 \xrightarrow{\sigma-1}R/\n^n\longrightarrow0
 \tag{10}
\end{equation}
is exact. At $n=1$ it becomes
\[
 0\longrightarrow k\longrightarrow\bar k
 \xrightarrow{x\mapsto x^q-x}\bar k\longrightarrow0,
\]
where $\bar k$ is the algebraic closure of $k$. For the induction step, use
\[
\begin{CD}
0@>>>R/\n@>>>R/\n^n@>>>R/\n^{n-1}@>>>0\\
@.@V{\sigma-1}VV@V{\sigma-1}VV@V{\sigma-1}VV@.\\
0@>>>R/\n@>>>R/\n^n@>>>R/\n^{n-1}@>>>0.
\end{CD}
\]
The snake lemma gives surjectivity on $R/\n^n$ and a kernel with $q^n$ elements.
The image of $A/\n^n$, already contained in that kernel, has that size.
Passing to the inverse limit yields the exact sequence
\[
 0\longrightarrow A\longrightarrow B\xrightarrow{\sigma-1}B\longrightarrow0
\]
by I.A.8. The author states that the proof for $B^\times$ is similar.
\end{proof}
% END SUPPORT LEMMA I.3.11

% BEGIN SUPPORT PROPOSITION I.3.10 PROOF
\section*{Author proof of Proposition I.3.10}
Composition inverses of series are written $h^{-1}$, so $h\circ h^{-1}=T=h^{-1}\circ h$.
The author divides the proof into four steps:
\begin{enumerate}
\item Construct $\theta\in B[[T]]$ satisfying (a),(b).
\item Choose it so that $g=\sigma\theta\circ f\circ\theta^{-1}$.
\item Show that $\theta(F_f(\theta^{-1}(X),\theta^{-1}(Y)))$ has the properties
characterizing $F_f(X,Y)$, hence equals it.\footnote{The original Step 3 says
$F_f$; the author's separate erratum for printed p.41 specifies $F_g$ here.
This source correction is identified rather than silently applied.}
\item Show that $\theta\circ[a]_f\circ\theta^{-1}$ has the properties characterizing
$[a]_g$, hence equals it.
\end{enumerate}

\subsection*{Step 1}
Choose $\varepsilon\in B^\times$ with $\sigma\varepsilon=\varepsilon u$ using I.3.11.
Starting from $\theta_1(T)=\varepsilon T$, construct polynomials with
\[
 \theta_r(T)=\theta_{r-1}(T)+bT^r\quad(b\in B),\qquad
 \sigma\theta_r=\theta_r\circ[u]_f+\terms{r+1}.
\]
The first case is $\sigma\varepsilon T=\varepsilon(uT+\cdots)+\terms{2}$.
Given $\theta_r$, set $\theta_{r+1}=\theta_r+bT^{r+1}$ and write
$b=a\varepsilon^{r+1}$. The desired relation is
\[
 (\sigma\theta_r)(T)+(\sigma a)(\sigma\varepsilon)^{r+1}T^{r+1}
 =\theta_r([u]_f(T))+a\varepsilon^{r+1}(uT)^{r+1}+\terms{r+1}.
\]
Thus the required coefficient equation is
\[
 (\sigma a-a)(\varepsilon u)^{r+1}=c,
\]
where $c$ is the coefficient of $T^{r+1}$ in
$\theta_r\circ[u]_f-\sigma\theta_r$.
Choose $a\in B$ with $\sigma a-a=c/(\varepsilon u)^{r+1}$.

\subsection*{Step 2}
Define
\[
 h=\sigma\theta\circ f\circ\theta^{-1}
 =\theta\circ[u]_f\circ f\circ\theta^{-1}
 =\theta\circ f\circ[u]_f\circ\theta^{-1}.
\]
Since $f$ and $[u]_f$ have coefficients in $A$,
\[
 \sigma h=\sigma\theta\circ f\circ[u]_f\circ\sigma\theta^{-1}
 =\sigma\theta\circ f\circ\theta^{-1}=h.
\]
Here $[u]_f\circ\sigma\theta^{-1}=\theta^{-1}$ follows from
$\theta\circ[u]_f\circ\sigma\theta^{-1}=T$.
It follows that $h\in A[[T]]$. Moreover,
\[
 h(T)=\sigma\varepsilon\cdot\pi\cdot\varepsilon^{-1}T+\cdots
 =\varpi T+\terms{2},
\]
and
\[
 h(T)\equiv\sigma\theta\circ(\theta^{-1})^q
 =\sigma\theta(\sigma\theta^{-1}(T^q))\equiv T^q\pmod{\m_K}.
\]
Thus $h\in\mathcal F_\varpi$. Set $\theta'=[1]_{g,h}\circ\theta$.
It continues to satisfy (a),(b), because $[1]_{g,h}\in A[[T]]$, and
\[
 \sigma\theta'\circ f\circ(\theta')^{-1}
 =[1]_{g,h}\circ h\circ[1]_{g,h}^{-1}=g.
\]
The author leaves Steps 3 and 4 as straightforward applications of I.2.11.
% END SUPPORT PROPOSITION I.3.10 PROOF

% BEGIN PRIMARY PROOF FIELD INDEPENDENCE
\section*{Author proof that $K_\pi K^{\un}$ is independent of $\pi$}
Let $\varpi=u\pi$. Proposition I.3.10 gives
\[
 (\sigma\theta)\circ[\pi]_f
 =\theta\circ[u]_f\circ[\pi]_f
 =\theta\circ[\varpi]_f=[\varpi]_g\circ\theta,
\]
or $(\sigma\theta)(f(T))=g(\theta(T))$.
For $\alpha\in K^{\al}$,
\[
 f(\alpha)=0\ \Longrightarrow\ g(\theta(\alpha))=0,
 \qquad
 g(\alpha)=0\ \Longrightarrow\ f(\theta^{-1}(\alpha))=0.
\]
Hence $\theta$ gives a bijection $\Lambda_{f,1}\longrightarrow\Lambda_{g,1}$, and
\begin{align*}
 \Khat[\Lambda_{g,1}]
 &=\Khat[\theta(\Lambda_{f,1})]
 \subset\Khat[\Lambda_{f,1}]\\
 &=\Khat[\theta^{-1}(\Lambda_{g,1})]
 \subset\Khat[\Lambda_{g,1}].
\end{align*}
Thus $\Khat[\Lambda_{g,1}]=\Khat[\Lambda_{f,1}]$.
The next lemma provides
\[
 \Khat[\Lambda_{g,1}]\cap K^{\al}=K^{\un}[\Lambda_{g,1}],\qquad
 \Khat[\Lambda_{f,1}]\cap K^{\al}=K^{\un}[\Lambda_{f,1}],
\]
so $K^{\un}[\Lambda_{g,1}]=K^{\un}[\Lambda_{f,1}]$.
The author extends the same argument to every $n$, obtaining
\[
 K^{\un}[\Lambda_{g,n}]=K^{\un}[\Lambda_{f,n}],\qquad
 K^{\un}\cdot K_\varpi=K^{\un}\cdot K_\pi.
\]
% END PRIMARY PROOF FIELD INDEPENDENCE

% BEGIN SUPPORT LEMMA I.3.12
\begin{closedfield}
Every subfield $E$ of $K^{\al}$ containing $K$ is closed in the topological sense.
\end{closedfield}
\begin{proof}[Author proof]
Put $H=\Gal(K^{\al}/E)$. It fixes $E$ and therefore, by continuity, every element
in the closure of $E$. Galois theory then identifies $E$ with its closure in $K^{\al}$.
\end{proof}
% END SUPPORT LEMMA I.3.12

% BEGIN PRIMARY PROOF ARTIN ACTION
\section*{Author proof that $\phi_\pi$ is independent of $\pi$}
For two prime elements $\pi,\varpi$, the author proves
$\phi_\pi(\varpi)=\phi_\varpi(\varpi)$.
Since $\pi$ is arbitrary, another prime element $\pi'$ then satisfies
\[
 \phi_{\pi'}(\varpi)=\phi_\varpi(\varpi)=\phi_\pi(\varpi).
\]
Prime elements generate $K^\times$, so this proves $\phi_\pi=\phi_{\pi'}$.
On $K^{\un}$ the two maps both induce Frobenius; compare them on $K_\varpi$.
The isomorphism $\theta:F_f\to F_g$ of I.3.10 induces
$\Lambda_{f,n}\to\Lambda_{g,n}$. By definition $\phi_\varpi(\varpi)$ is the identity
on $K_\varpi$. Since the elements $\theta(\lambda)$ for $\lambda\in\Lambda_{f,n}$
generate $K_{\varpi,n}$, it remains to show
\[
 \phi_\pi(\varpi)(\theta(\lambda))=\theta(\lambda),\qquad
 \lambda\in\Lambda_{f,n}.
\]
Write $\varpi=u\pi$ and
$\phi_\pi(\varpi)=\phi_\pi(u)\phi_\pi(\pi)=\tau\sigma$, where
\[
 \sigma=\begin{cases}\Frob_K&\text{on }K^{\un},\\\mathrm{id}&\text{on }\lambda,\end{cases}
 \qquad
 \tau=\begin{cases}\mathrm{id}&\text{on }K^{\un},\\{}[u^{-1}]_f&\text{on }\lambda.\end{cases}
\]
Using the coefficients of $\theta$ in $\Khat$ and I.3.10, the author's final identity is
\[
 \phi_\pi(\varpi)(\theta(\lambda))
 =\tau\sigma(\theta(\lambda))
 =(\sigma\theta)(\tau\lambda)
 =(\sigma\theta)([u^{-1}]_f(\lambda))
 =\theta(\lambda).
\]
% END PRIMARY PROOF ARTIN ACTION

\section*{Source and revision map}
The primary statement, comparison proposition, difference lemma and Artin-map definition
are on physical p.49 / printed p.40. The difference-lemma proof and start of Step 1 are on
p.50 / p.41; Step 1, Step 2 and the first field-comparison argument continue on
p.51 / p.42. The completion-to-algebraic comparison, closed-field lemma and final Artin
comparison are on p.52 / p.43. Lemma I.2.11 is on p.41--42 / p.32--33.
The formal-module summary is p.44 / p.35; torsion definitions are p.45 / p.36;
the structure results are p.47 / p.38.
The complete author PDF and the author's separately saved errata remain the authoritative
sources. No support result in this reader is counted as another corpus problem.
\end{document}
